\documentclass[11pt]{article}
\textwidth 15.9cm
\textheight 23. cm
\addtolength{\oddsidemargin}{-11.5mm}
\addtolength{\evensidemargin}{-11.5mm}
\addtolength{\topmargin}{-19.0mm}

%\usepackage{amsmath}
\usepackage[T1]{fontenc}
\usepackage[ansinew]{inputenc}
\usepackage[bbgreekl]{mathbbol}
\usepackage{amsmath,amssymb}
\usepackage{graphicx}
\usepackage[colorlinks=true,urlcolor=cyan]{hyperref}%,linkcolor=green]{hyperref}
\usepackage{url}
\usepackage[numbers,square]{natbib}
\usepackage{multirow}
\usepackage{tikz}
\usepackage{tikz-cd}
\usetikzlibrary{arrows,matrix}
\usepackage{bbm}
\usepackage{authblk}
\usepackage{subcaption}
\usepackage{xstring}
\usepackage{xparse}
\usepackage{comment}
\renewcommand{\Affilfont}{\footnotesize}

\newcommand{\amrex}{\texttt{AMReX}}

\usepackage{stmaryrd}
\usepackage{comment}
\usepackage{showlabels}
\usepackage{soul}

\usetikzlibrary{calc,patterns,shapes.geometric}
\usetikzlibrary{arrows,shapes,positioning}
\usetikzlibrary{decorations.markings,decorations.pathmorphing,decorations.pathreplacing}
\usetikzlibrary{decorations.text,decorations.shapes}

%%%%%%%% MACROS
\input{GempicLatexMacros}

\lstdefinelanguage{pkgconfig}{
     morecomment=[l]{\#}
}

\title{Input Files}


\author[1]{The Gempic development team}
\affil[1]{Max-Planck-Institut f\"ur Plasmaphysik, Garching, Germany}


\begin{document}

\maketitle

\begin{abstract}
This note describes the input file semantics and their meaning
\end{abstract}

\tableofcontents

\section{Input Files}
The \amrex{} framework provides a simple way to give input parameters to a simulation: The input file. This is a file in ordinary \verb|.txt| encoding (typically named \verb|.input|) that can be read as an argument to an executable compiled with \amrex:
\begin{lstlisting}[language=sh]
    $ ./executable inputFile.input 
\end{lstlisting}

\subsection{The Broad Strokes of Input Files}
An input file consists of separate lines of variables, each with the variable name, an equals (\verb|=|) sign and one or more values. Variables can have an optional scope, typically corresponding to a class, and values are separated by spaces:
\begin{lstlisting}[language=pkgconfig]
# Comments begin with '#'
# Input file for simulation etc
simName = "etc"

# Input for class Filter
Filter.enable = 1 # enable filtering
Filter.nPass = 3 3 3 # pass the filter three times in each direction
\end{lstlisting}
%
Notice that, like the class and variable counterparts:
\begin{itemize}
     \item Scope names are \verb|CamelCase|
     \item Variable names are \verb|camelBack|
\end{itemize}
%
In the following, the possible options are discussed for each category of input variables. Most of these will have default options that are used if the values are not provided by the user.

\subsection{Global Variables}
A list of global variables (\emph{i.e.}, variables read in more than one place) follows below. Note that these variables \emph{never} have a scope. A fully up-to-date lists can always be found in \verb|GEMPIC_Parameters.H|.
\begin{lstlisting}[language=pkgconfig]
# Global variables
# The name of the output file.
outputFile = "outputFileName" # String. No default.

# The name of the simulation itself
simName = "my_simulation" # String. No default.

# How many additional ghost cells to put. Note that Gempic always calculates a default
# number of ghost cells in addition to this, depending on Hodge- and spline parameters.
nGhostExtra = 2 # Integer. Default 3.

# Wave vector. This determines the physical length of the mesh, which is equal to 2pi/k.
# If this is given, ComputationalDomain.domainHi will therefore be ignored.
k = 1 1 0.3 # Real vector of AMReX_SPACEDIM entries. No default.
\end{lstlisting}

\subsection{ComputationalDomain}
The computational domain is the region of space in which the simulation is performed. The domain is divided into a number of cells, given by the \verb|nCell| variable. The mesh is then divided into grids, each with a maximum size given by the \verb|maxGridSize| variable.

\begin{lstlisting}[language=pkgconfig]
# Computational domain
# Number of cells used in the grid
ComputationalDomain.nCell = 32 8 8 # Integer vector of AMReX_SPACEDIM entries. No default.

# Maximum grid size for partitioning computations. 
# The maxGridSize must be less than or equal to nCell,
# and in each direction, nCell must be divisible by maxGridSize.
# Here we would have 4 regions in the x-direction and 1 in the y- and z-directions.
ComputationalDomain.maxGridSize = 8 8 8 # Integer vector of AMReX_SPACEDIM entries. No default.

# Is this direction periodic?
ComputationalDomain.isPeriodic = 1 1 1 # Boolean (1/0) vector of AMReX_SPACEDIM entries. No default.

# Lower physical coordinates of the mesh
ComputationalDomain.domainLo = 1.0 0.0 0.0 # Real vector of AMReX_SPACEDIM entries. Default 0.0 0.0 0.0.

# Higher physical coordinates of the mesh. Will be ignored if the wave vector k is given.
ComputationalDomain.domainHi = 10.0 10.0 5.0 # Real vector of AMReX_SPACEDIM entries. Default 1.0 1.0 1.0.

\subsection{Filter}\label{sec:FilterInput}
Controls the filtering behaviour of the simulation (assuming your simulation has this option).
See \ref{sec:BilinearFilter} for details.
\begin{lstlisting}[language=pkgconfig]
# Filter class options

# To filter or not to filter 
Filter.enable = 1 # Boolean (1/0). Default 0.

# How many times to pass the filter in each direction
Filter.nPass = 2 2 2 # Integer vector of AMReX_SPACEDIM entries. No default.

# Whether or not to do an extra compensation step when calculating the filter stencil
Filter.compensate = 1 # Boolean (1/0). Default 0.
\end{lstlisting}

\subsection{FullDiagnostics}
Controls the behaviour of the full diagnostics module (assuming your simulation has this option).
\begin{lstlisting}[language=pkgconfig]
# FullDiagnostics class options
# Names of individual components of fields that have several components.
# Note that components are not the same as directions.
# Note that the names are not explicitly given inside the FullDiagnostics scope.
componentNames.{fieldName} = "density" "total pressure" # String vector. Default "var0" "var1" ...

FullDiagnostics.enable = 1 # Boolean (1/0). Default 0.

# Arbitrary names for groups of diagnostics. No naming style is enforced.
FullDiagnostics.groupNames = any group_name yOu want # String vector. No default.

# How often to output diagnostics. Overwrites default and is used for each group 
# unless they specify values themselves.
FullDiagnostics.saveInterval = 55 # Integer. No default.

# Where to save the output files. Overwrites default and is used for each group
# unless they specify values themselves.
FullDiagnostics.saveFolder = "FD/1/" # String. Default "FullDiagnostics/".

## Options for each individual group
# Where to save the output files for this specific group
FullDiagnostics.{groupName}.saveFolder = "FD/2/" # String. Default "FullDiagnostics/".

# Prefix to output files
FullDiagnostics.{groupName}.filePrefix = "." # String. Default "plt_{groupName}".

# Minimal number of digits of output file names, i.e. rho000001
FullDiagnostics.{groupName}.fileMinDigits = 5 # Integer. Default 6.

# If true, a dump is performed at the last timestep regardless of the required dump timesteps
FullDiagnostics.{groupName}.dumpLastTimeStep = 0 # Boolean (1/0). Default 1.

# Names of the variables in this group to output. Must correspond to names in simulation.
# Cannot be a mixture of particles and fields.
FullDiagnostics.{groupName}.varNames = Ey Ez Bx # String vector. No default.

# Format of the output files. Options are "plotfile" (and maybe some day "openPMD")
FullDiagnostics.{groupName}.format = "plotfile" # String. Default "plotfile".

# Output processor. Only applicable to fields. Options are "Raw" and "CellCenter", determining
# if you want the raw data or the interpolated values of the fields in the cell centers.
FullDiagnostics.{groupName}.outputProcessor = "Raw" # String. Default "Raw".

# How often to output the diagnostics. Every n time steps.
FullDiagnostics.{groupName}.saveInterval = 10 # Integer. Default 100.

\end{lstlisting}

\subsection{ReducedDiagnostics}
Controls the behaviour of the reduced diagnostics module (assuming your simulation has this option).
\begin{lstlisting}[language=pkgconfig]
# ReducedDiagnostics class options

ReducedDiagnostics.enable = 1 # Boolean (1/0). Default 0.

# Where or not to compute fields that might be needed for reduced calculations
# even if they aren't initiated in the main script.
ReducedDiagnostics.computeMissingFields = 1 # Boolean (1/0). Default 0.

# Arbitrary names for groups of diagnostics. No naming style is enforced.
ReducedDiagnostics.groupNames = any group_name yOu want # String vector. No default.

# How often to output diagnostics. Overwrites default and is used for each group 
# unless they specify values themselves.
ReducedDiagnostics.saveInterval = 55 # Integer. No default.

# Where to save the output files. Overwrites default and is used for each group
# unless they specify values themselves.
ReducedDiagnostics.saveFolder = "RD/1/" # String. Default "ReducedDiagnostics/".

# File extension of the output files. Overwrites default and is used for each group
# unless they specify values themselves.
ReducedDiagnostics.extension = "output" # String. Default "txt".

## Options for each individual group
# Names of the variables in this group to output. Options are:
# Particle ElecFieldEnergy MagFieldEnergy GaussError
ReducedDiagnostics.{groupName}.types = Particle GaussError # String vector. No default.

# How often to output the diagnostics. Every n time steps.
ReducedDiagnostics.{groupName}.saveInterval = 10 # Integer. Default 1.

# Where to save the output file of this specific group
ReducedDiagnostics.{groupName}.saveFolder = "RD/2/" # String. Default "ReducedDiagnostics/".

# Name of output file of this specific group. Can include extension, which will be
# overwritten by the .extension parameter if given.
ReducedDiagnostics.{groupName}.fileName = lottery_numbers # String. Default {groupName}.

# File extension of the output file of this specific group
ReducedDiagnostics.{groupName}.extension = "output_special" # String. Default "txt".

\end{lstlisting}

\subsection{Parameters}
The Parameters class itself has the setting to print a new input file with all of the parameters used in the simulation.\\
For this to have any effect, however, one must call \lstinline|Parameters::set_print_output(true);| in the simulation \emph{before the first Parameters object is initialized}.
\begin{lstlisting}[language=pkgconfig]
# Parameters parameters
outputFile = "simulationParameters.input"
\end{lstlisting}

\subsection{Particle}\label{sec:ParticleInput}
Particle input variables are divided into species with subscopes given by user provided \verb|speciesNames|:
\begin{lstlisting}[language=pkgconfig]
# Particle parameters
# Note that no naming style is enforced.
# The name has to match that given to Reduced- and FullDiagnostics.
Particle.speciesNames = any_names yOU WANT # String vector. No default.
Particle.sampler = PseudoRandom # String. Default is PseudoRandom.
                                # Can be PseudoRandom|PseudoRandomBox|Sobol|Rejection|RejectionSobol
Particle.seed = 1234 # Integer. Default is 925.
# Number of particles per cell
# Parameters for the first species
Particle.{speciesName}.sampler = Sobol # overrides general sampling scheme for {speciesName}.
Particle.{speciesName}.nPartPerCell = 10 # Integer. No default.
Particle.{speciesName}.charge = -1.0 # Real. Default 1.0.
Particle.{speciesName}.mass = 1.0 # Real. Default 1.0.
Particle.{speciesName}.compNames =  "vx" "vy" "vz" "weight" # List of strings. Default. Names extra
                                                            # components in e.g. output files.

# Initial density of the particle species using the Parser class in AMReX. Here, kvarx=k[0].
Particle.{speciesName}.density = "1.0 + 0.01 * cos(kvarx * x)" # String. No default.

# Number of Gaussians making up the velocity distribution of the {speciesName}
Particle.{speciesName}.numGaussians = 2 # Integer. Default 1.

# Gaussian parameters for the first (G0) Gaussian. The sum of all Gaussian weights must be 1.0.
Particle.{speciesName}.G0.vMean = 0.0 0.0 0.0 # Real vector. No default.
Particle.{speciesName}.G0.vThermal = 1 1 1 # Real vector. No default.
Particle.{speciesName}.G0.vWeight = 0.7 # Real. No default.

# Gaussian parameters for the second (G1) Gaussian. The sum of all Gaussian weights must be 1.0.
# When using functions to specify particle velocities, all directions must be supplied.
Particle.{speciesName}.G1.vMean = 0.0 0.0 0.0 # Real vector. No default.
Particle.{speciesName}.G1.vThermal.x = "cos(kvarx * x)" # String. No default.
Particle.{speciesName}.G1.vThermal.y = "1.0" # String. No default.
Particle.{speciesName}.G1.vThermal.z = "1.0" # String. No default.
Particle.{speciesName}.G1.vWeight = 0.3 # Real. No default.
\end{lstlisting}
See \ref{sec:ParticleStructure} for details.

For rejection sampling, the \verb|density| and \verb|G0...| configure the proposal distribution. The additional inputs are
\begin{lstlisting}[language=pkgconfig]
Particle.{speciesName}.sampler = Rejection
Particle.{speciesName}.rejectionDistribution = "1 / sqrt(2*3.14159265)^3 * exp(-(vx^2 + vy^2 + vz^2)/2)" # String. No default.
Particle.{speciesName}.rejectionScaling = 2.0 # Real. Default 1.0.
Particle.{speciesName}.nMaxRejections = 100 # Integer. Default 100.
\end{lstlisting}

In case of a one-species simulation, it is sometimes necessary to add a background to rho to guarantee a neutral plasma. This can be achieved via the input parameter 

\begin{lstlisting}
rhoBackground = 1
\end{lstlisting}

If this variable is not set, the background is automatically assumed to be 0. 
An example for parts of a single-species input file is

\begin{lstlisting}[language=pkgconfig]
# Particle parameters
Particle.speciesNames = electrons
rhoBackground = 1 # one species

Particle.electrons.nPartPerCell = 1000 # number of particles per cell
Particle.electrons.charge = -1.0 # vector (one entry per species)
Particle.electrons.mass = 1.0 # vector (one entry per species)
Particle.electrons.density = "1.0 + 0.04 * cos(kvarx * x)"
Particle.electrons.numGaussians = 1
\end{lstlisting}

while an example for parts of a two-species input file is

\begin{lstlisting}[language=pkgconfig]
# Particle parameters
Particle.speciesNames = electrons ions
rhoBackground = 0 # not necessary, default value in simulation file
    
Particle.electrons.nPartPerCell = 1000 # number of electrons per cell
Particle.electrons.charge = -1.0 # vector (one entry per species)
Particle.electrons.mass = 1.0 # vector (one entry per species)
Particle.electrons.density = "1 + 0.04 * cos(kvarx * x)"
Particle.electrons.numGaussians = 1
    
Particle.ions.nPartPerCell = 1000 # number of ions per cell 
Particle.ions.charge = 1.0 # vector (one entry per species)
Particle.ions.mass = 100.0 # vector (one entry per species)
Particle.ions.density = "1.0" 
Particle.ions.numGaussians = 1
\end{lstlisting}

\subsection{Function}
Variables read in the simulation by the helper function interpreter \verb|parse_funtions| using the \href{https://amrex-codes.github.io/amrex/docs_html/Basics.html#parser}{Parser class in AMReX}. The variable names can be anything.
\begin{lstlisting}[language=pkgconfig]
# Functions
Function.Bz = "1e-3 * cos(kvarx * x)" # String. No default.
\end{lstlisting}

\subsection{BoundaryCondition}
Controls the boundary conditions on non-periodic domains (assuming your simulation supports this option). If the simulation domain contains non-periodic boundaries, i.e., \verb|isPeriodic| contains one or more 0 entries, boundary conditions can be provided. Specific boundary conditions can only be assigned to named fields. For all other fields the default is used.
Particle boundary conditions must be implemented separately.
\begin{lstlisting}[language=pkgconfig]
# Boundary conditions for fields
BoundaryCondition.E = PerfectlyConducting PerfectlyConducting # String vector. 2 entries for each non-periodic dimension (lower/upper)
BoundaryCondition.Default = Natural Natural # String vector. Default "Natural" * 2 * dim_non_periodic
\end{lstlisting}

\subsection{PoissonSolver}\label{sec:PoissonSolverInput}
Selects which solver to use for the poisson equation (\ref{eq:Poisson_example}), provided the PoissonSolver is used in the simulation.
There are several choices built in:
\begin{itemize}
	\item \verb|FFT|: Using the Fast Fourier Transform to solve the Poisson equation in Fourier space. Default if the boundary conditions are periodic and FFT was found during installation. Does not use \verb|relTol| and \verb|absTol|.
	\item \verb|Hypre|: Using the sparse solver library Hypre. Default if Hypre is available and FFT is not applicable.
	\item \verb|ConjugateGradient|: Uses our own conjugate gradient solver. Default if neither FFT nor Hypre apply.
	\item \verb|ConjugateGradientInverseHodge|: Same as ConjugateGradient, but using another conjugate gradient solver instead of the Hodge operator such that $-\tildemat{D}{\mat{H}}_2\mat{G}\arr{\phi}=\arr{\rho}$ exactly.
	\item \verb|Amrex|: Using the AMReX MLMG solver. Boundary conditions must be periodic.
\end{itemize}
\begin{lstlisting}[language=pkgconfig]
# PoissonSolver
PoissonSolver.solver = ConjugateGradient # String. Default: See above.
PoissonSolver.relTol = 1e-11 # Real. Default 1e-11.
PoissonSolver.absTol = 1e-12 # Real. Default 1e-12.
\end{lstlisting}

\subsection{HYPRE options}
These parameters are for the \verb|HYPRE| library used to parallely solve large sparse linear systems. These can be found in the \verb|AMReX_HypreIJIface.cpp| file.

There are 8 solver choices available:
\begin{lstlisting}[language=pkgconfig]
# hypre solver
hypre.hypre_solver = GMRES # BoomerAMG(Default) GMRES COGMRES LGMRES FlexGMRES BiCGSTAB PCG Hybrid
\end{lstlisting}

There are 3 preconditioner choices available: 
\begin{lstlisting}[language=pkgconfig]
# hypre preconditioner
hypre.hypre_preconditioner = BoomerAMG # BoomerAMG euclid none(Default. the default solver option "hypre.hypre_solver=BoomerAMG" will only work with "hypre.hypre_preconditioner=none")
\end{lstlisting}

Parameters for the \verb|BoomerAMG| preconditioner:
\begin{lstlisting}[language=pkgconfig]
# hypre BoomerAMG preconditioner configure
hypre.bamg_max_iterations = 20 # integer. Default:20
hypre.bamg_max_levels = 25 # integer. Default:25
hypre.bamg_smooth_type = 6 # integer. Default:6
\end{lstlisting}

Depending on the value of the \verb|hypre.bamg_smooth_type| paramater used for the \verb|BoomerAMG| preconditioner, various other sub-parameters can be configured.
Sub-parameters for \verb|hypre.bamg_smooth_type = 5|:
\begin{lstlisting}[language=pkgconfig]
# hypre options for BoomerAMG preconditioner for hypre.bamg_smooth_type = 5
hypre.bamg_smooth_num_sweeps = 1 # integer. Default:1
hypre.bamg_smooth_num_levels = 0 # integer. Default:0
hypre.bamg_ilu_type = 0 # integer. Default:0
hypre.bamg_ilu_level = 0 # integer. Default:0
hypre.bamg_ilu_max_iter = 1 # integer. Default:1
hypre.bamg_ilu_reordering_type = 1 # integer. Default:1 (1==RCM, 0==None)
hypre.bamg_ilu_tri_solve = # integer. Default:1 0-iterative, 1-direct
hypre.bamg_ilu_upper_jacobi_iters = 5 # integer. Default:5
hypre.bamg_ilu_lower_jacobi_iters = 5 # integer. Default:5
\end{lstlisting}

Sub-parameters for \verb|hypre.bamg_smooth_type = 7|:
\begin{lstlisting}[language=pkgconfig]
# hypre options for BoomerAMG preconditioner for hypre.bamg_smooth_type = 7
hypre.bamg_smooth_num_sweeps = 1 # integer. Default:1
hypre.bamg_smooth_num_levels = 0 # integer. Default:0
hypre.bamg_ilu_max_iter = 1 # integer. Default:1
hypre.bamg_ilu_max_row_nnz = 1000 # integer. Default:1000
hypre.bamg_ilu_drop_tol = 1.e-10 # real. Default:(1.0e-02 for hypre.bamg_smooth_type = 7 and 1.e-10 otherwise)
\end{lstlisting}

Sub-parameters for \verb|hypre.bamg_smooth_type = 9|:
\begin{lstlisting}[language=pkgconfig]
# hypre options for BoomerAMG preconditioner for hypre.bamg_smooth_type = 9
hypre.bamg_smooth_num_sweeps = 1 # integer. Default:1
hypre.bamg_smooth_num_levels = 0 # integer. Default:0
\end{lstlisting}

Parameters for the \verb|Euclid| preconditioner:
\begin{lstlisting}[language=pkgconfig]
# hypre Euclid preconditioner configure
hypre.euclid_level = 4 # integer. Default:-1. for 3D problems set to 1, for 2D set to 4-8
hypre.euclid_use_block_jacobi = 0 # integer. Default:0 (0==PILU, 1==Block Jacobi)
hypre.euclid_stats = 0 # integer. Default:0 (0 1)
hypre.euclid_mem = 0 # integer. Default:0 (0 1)
\end{lstlisting}

Parameters for the \verb|BoomerAMG| solver:
\begin{lstlisting}[language=pkgconfig]
# hypre BoomerAMG solver configure
hypre.verbose = 0 # integer. Default:0. 1(print setup info) 2(print solve info) 3(print setup+solve info)
hypre.logging = 0 # integer. Default:0. >1(latest residual available)
\end{lstlisting}

Parameters for the \verb|GMRES| solver:
\begin{lstlisting}[language=pkgconfig]
# hypre GMRES solver configure
hypre.verbose = 0 # integer. Default:0. 1(print setup info) 2(print solve info) 3(print setup+solve info)
hypre.logging = 0 # integer. Default:0. >1(latest residual available)
hypre.num_krylov = 50 # integer. Default:50
hypre.max_iterations = 200 # integer. Default:200
\end{lstlisting}

Parameters for the \verb|COGMRES| solver:
\begin{lstlisting}[language=pkgconfig]
# hypre COGMRES solver configure
hypre.verbose = 0 # integer. Default:0. 1(print setup info) 2(print solve info) 3(print setup+solve info)
hypre.logging = 0 # integer. Default:0. >1(latest residual available)
hypre.num_krylov = 50 # integer. Default:50
hypre.max_iterations = 200 # integer. Default:200
\end{lstlisting}

Parameters for the \verb|LGMRES| solver:
\begin{lstlisting}[language=pkgconfig]
# hypre LGMRES solver configure
hypre.verbose = 0 # integer. Default:0. 1(print setup info) 2(print solve info) 3(print setup+solve info)
hypre.logging = 0 # integer. Default:0. >1(latest residual available)
hypre.num_krylov = 50 # integer. Default:50
hypre.max_iterations = 200 # integer. Default:200
\end{lstlisting}

Parameters for the \verb|FlexGMRES| solver:
\begin{lstlisting}[language=pkgconfig]
# hypre FlexGMRES solver configure
hypre.verbose = 0 # integer. Default:0. 1(print setup info) 2(print solve info) 3(print setup+solve info)
hypre.logging = 0 # integer. Default:0. >1(latest residual available)
hypre.num_krylov = 50 # integer. Default:50
hypre.max_iterations = 200 # integer. Default:200
\end{lstlisting}

Parameters for the \verb|BiCGSTAB| solver:
\begin{lstlisting}[language=pkgconfig]
# hypre BiCGSTAB solver configure
hypre.verbose = 0 # integer. Default:0. 1(print setup info) 2(print solve info) 3(print setup+solve info)
hypre.logging = 0 # integer. Default:0. >1(latest residual available)
hypre.max_iterations = 200 # integer. Default:200
\end{lstlisting}

Parameters for the \verb|PCG| solver:
\begin{lstlisting}[language=pkgconfig]
# hypre PCG solver configure
hypre.verbose = 0 # integer. Default:0. 1(print setup info) 2(print solve info) 3(print setup+solve info)
hypre.logging = 0 # integer. Default:0. >1(latest residual available)
hypre.max_iterations = 200 # integer. Default:200
\end{lstlisting}

Parameters for the \verb|Hybrid| solver:
\begin{lstlisting}[language=pkgconfig]
# hypre Hybrid solver configure
hypre.verbose = 0 # integer. Default:0. 1(print setup info) 2(print solve info) 3(print setup+solve info)
hypre.logging = 0 # integer. Default:0. >1(latest residual available)
hypre.num_krylov = 50 # 50(default)
hypre.hybrid_dscg_max_iter = 1000 # integer. Default:1000
hypre.hybrid_pcg_max_iter = 100 # integer. Default:100
hypre.hybrid_setup_type = 1 # integer. Default:1
hypre.hybrid_solver_type = 1 # integer. Default:1(PCG) 2(GMRES) 3(BiCGSTAB)
\end{lstlisting}

\subsection{Example of An Input File}
Putting it all together, an input file might look like this:
\begin{lstlisting}[language=pkgconfig]
# Input file for adiabatic slab simulation
simName = "test_adiabatic_slab"

# output file with read and computed parameters
outputFile = "test_adiabatic_slab_parameters.input"

# Grid parameters
ComputationalDomain.nCell = 128 8 8
ComputationalDomain.maxGridSize = 32 8 8
ComputationalDomain.isPeriodic = 1 1 1 # 1 -> periodic, 0 -> else
k = 1 1 0.3

# Particle parameters
Particle.speciesNames = electrons ions

# Electrons
Particle.electrons.nPartPerCell = 10
Particle.electrons.charge = -1.0
Particle.electrons.mass = 1.0
Particle.electrons.density = "1.0 +  cos(kvarx * x) + sin(kvary * y) + cos(2*kvarz *z)"

Particle.electrons.numGaussians = 2
Particle.electrons.G0.vMean = 0 0 0
Particle.electrons.G0.vThermal = 1.0 1.0 1.0
Particle.electrons.G0.vWeight = 0.75
Particle.electrons.G1.vMean = -2.4 0 0
Particle.electrons.G1.vThermal = 1.0 1.0 1.0
Particle.electrons.G1.vWeight = 0.25

# Ions
Particle.ions.nPartPerCell = 100 
Particle.ions.charge = 1.0
Particle.ions.mass = 1.0
Particle.ions.density = "1.0 +  0.2*cos(kvarx * x)"

Particle.ions.numGaussians = 1
Particle.ions.G0.vMean = 2.4 0 0
Particle.ions.G0.vThermal = 1.0 1.0 1.0
Particle.ions.G0.vWeight = 1.0

# Filter
Filter.enable = 1
Filter.nPass = 3 3 3
Filter.compensate = 1

# Full Diagnostics
FullDiagnostics.enable = 1
FullDiagnostics.groupNames = part_test Ez_raw plots 

FullDiagnostics.part_test.varNames = ions

FullDiagnostics.plots.varNames = Ey Ez Bx
FullDiagnostics.plots.outputProcessor = CellCenter

FullDiagnostics.Ez_raw.varNames = Ez
FullDiagnostics.Ez_raw.saveInterval = 50
FullDiagnostics.Ez_raw.outputProcessor = Raw

# Reduced Diagnostics 
ReducedDiagnostics.groupNames = PartDiag FieldDiag GaussErr 
ReducedDiagnostics.saveInterval = 1 # Default for all reduced diagnostic groups

ReducedDiagnostics.FieldDiag.types = ElecFieldEnergy MagFieldEnergy 

ReducedDiagnostics.PartDiag.types = Particle 
ReducedDiagnostics.PartDiag.saveInterval = 10 # Special for this diagnostic group

ReducedDiagnostics.GaussErr.types = GaussError

# Functions
Function.Bz = "1e-3 * cos(kvarx * x)"
\end{lstlisting}

\end{document}

